\section{Actual root intervals, compatible covers and modular constraints}
\label{sec:quartic-modular-checkpoint}

The sixth continuation advances three distinct arithmetic interfaces.
Counting actual integer roots in an interval, with their exact rank
stratification, removes an exponent factor from the prime-window error.
This proves a larger finite-window result for most actual roots while
retaining both the exceptional roots and the primes outside that window.
A separate descent through a class-number-one quartic field preserves
the conjugacy of the four factors. It reduces the actual coefficient
choices and gives explicit smooth complete intersections over $\mathbb Q$.
The uniform height of their individual points remains uncontrolled.

The independent modular direction attaches an actual degree-two Frey
$\mathbb Q$-curve to the second norm. Complete ordinary local calculations
and explicitly named representation-theoretic inputs restrict the level
and weight of a necessary residual newform. The candidate dimension budget
includes both finite-flat and non-finite-flat branches at the exponent
prime. These spaces need not be empty. An exactly studied non-CM boundary
curve explains why one specified finite residue-product elimination can
fail; its positive primitive local shadows are not global power solutions.

Ordinary proofs and independent agent review preceded eighteen new scoped
Lean declarations verifying the actual Frey polynomial arithmetic. The
full counting, geometric, Tate-algorithm and modularity arguments retain
their ordinary status and external dependencies. No unrestricted signed
tail estimate, uniform point-height theorem, or unbounded two-step radical
family is supplied. Standard $abc$ remains unproved and undisproved.
